% % Chapter 5 Results and Evaluation   (we have some choices here)

% % Illustration the system – few sample run throughs with sample queries and recommendations -tying back to the design and the various components. 
% % This will show the reader how the system works.

% % Evaluation

% % Talk about the classic metrics ( can cover in Chap 2) and they limitations in complex system. 
% % User -evaluation – tasks for users to undertake. Questions asked – we need to design these; Objective measure of the recommendations; Subjective feedback.

% % Efficiency issues





\section{Recommendation System Evaluation}
\label{sec:recommendation_system_evaluation}
The recommendation system fulfils its purpose of providing a proof of concept for a recommendation system that is adjustable and explainable.

\subsection{Sample Runs}
\label{sec:sample_runs}

\subsubsection{User Preference Extraction}
\label{sec:user_preference_extraction}
The system is able to extract user preferences from the conversation and use them to provide recommendations.
For example, in the following conversation, the user is able to provide their preferences and the system is able to extract them and use them to provide recommendations:



\section{Large Language Model Evaluation}
\label{sec:large_language_model_evaluation}
The large language model did as it was told as it was a good boy and did not misbehave.


\section{User Evaluation Results}
\label{sec:user_evaluation_results}


\subsection{Survey Results}
\label{sec:survey_results}
You can find the full survey results in \cref{appendix:surveyResults}.
Overall, the survey results were really positive, with the users finding the system quite useful, able to understand their preferences, and has a lot of potential.


\subsection{User Feedback}


\paragraph{Suggestions for Improvement:} Users provided valuable feedback for enhancing the system. 
They suggested including posters of the recommended films and providing a summary or preliminary background information of the profiles of other users with similar preferences, that are used to generate the collaborative filtering based recommendations.



